\section{An exact three-phase determinant and joint finite capacity}
\label{jpb-section}

The joint packet argument has an arithmetic refinement that keeps all
three phases inside one actual additive relation. It improves a constant
in the determinant comparison, without providing the missing global
radical estimate.

\begin{lemma}[The three actual phase determinants]\label{jpb-identity}
For integer Eisenstein coordinate pairs $z,w$, put
\[
 N(a,b)=a^2+ab+b^2,\qquad R(a,b)=(-b,a+b),
\]
and define
\[
 x=\det(z,w),\quad y=\det(z,R^2w),\quad d=\det(z,Rw).
\]
Then
\begin{align}
 d&=x+y, & N(z)N(w)&=x^2+xy+y^2,\label{jpb-norm}\\
 4\bigl(N(z)N(w)\bigr)^3-27(xyd)^2
   &=\bigl((x-y)(2x+y)(x+2y)\bigr)^2.\label{jpb-discriminant}
\end{align}
If an integer $m$ divides $xyd$ and
\begin{equation}\label{jpb-size}
 4\bigl(N(z)N(w)\bigr)^3<27m^2,
\end{equation}
then at least one of $x,y,d$ is zero.
\end{lemma}
\begin{proof}
Writing $z=(a,b)$ and $w=(c,e)$ gives
$x=ae-bc$, $y=ac+bc+be$ and $d=ac+ae+be$.
These expressions give \eqref{jpb-norm}. Expanding
$4(x^2+xy+y^2)^3-27x^2y^2(x+y)^2$ proves
\eqref{jpb-discriminant}. In particular
$27(xyd)^2\le4(N(z)N(w))^3<27m^2$.
But $m^2\mid(xyd)^2$, so this nonnegative integer strictly below
$m^2$ must vanish. The strict premise excludes $m=0$ automatically;
negative moduli and zero coordinate pairs require no extra convention.
\end{proof}

\begin{proposition}[Different phases, one product modulus]\label{jpb-moduli}
Let $(m_i)_{i\in S}$ be finitely many pairwise coprime integer moduli.
If each $m_i$ divides at least one of $x,y,d$, then
$\prod_{i\in S}m_i\mid xyd$. If \eqref{jpb-size} holds for this
product modulus, one phase determinant vanishes.
\end{proposition}
\begin{proof}
Each modulus divides the full product $xyd$. Pairwise coprimality
implies divisibility by their product, and Lemma~\ref{jpb-identity}
applies. The selected phase may differ for every $i$.
\end{proof}

\begin{proposition}[A finite joint-capacity theorem]\label{jpb-capacity}
Let $A$ and $S$ be finite sets, with $A$ represented by actual integer
pairs $z_u$, and let $f_i:A\to H_i$ be maps to finite sets for $i\in S$.
Suppose:
\begin{enumerate}
\item the integer moduli $m_i$ are pairwise coprime;
\item $f_i(u)=f_i(v)$ forces $m_i$ to divide one of the three phase
determinants for $z_u,z_v$;
\item for every $u,v$, \eqref{jpb-size} holds with $z=z_u$, $w=z_v$
and $m=\prod_i m_i$;
\item a zero phase determinant for $z_u,z_v$ implies $u=v$.
\end{enumerate}
Then the joint map $u\mapsto(f_i(u))_{i\in S}$ is injective, and
\[
 |A|\le\prod_{i\in S}|H_i|\le n^{|S|}
 \quad\text{if }n\in\mathbb N\text{ and }|H_i|\le n\text{ for every }i.
\]
\end{proposition}
\begin{proof}
Equality of the joint images gives a phase divisibility at every
coordinate. Proposition~\ref{jpb-moduli} and the size bound force a
zero determinant, so the separation premise identifies the two inputs.
Counting the finite product target proves both inequalities.
\end{proof}

\begin{lemma}[Private independence survives a cubic kernel]
\label{jpb-kernel}
Let $G$ be a commutative group, $\phi:G\to H$ a group homomorphism
whose kernel is killed by three, and $(a_i)_{i\in I}$ a finite family
with integer-valued homomorphisms $v_i$ satisfying
$v_i(a_j)=0$ for $j\ne i$ and $v_i(a_i)\ne0$.
Then $x\mapsto\phi(\prod_i a_i^{x_i})$ is injective on $\mathbb Z^I$.
\end{lemma}
\begin{proof}
Equality of two images places their quotient in the kernel. Cubing
and applying $v_i$ gives $3v_i(a_i)(x_i-y_i)=0$. The nonzero diagonal
value gives $x_i=y_i$ for every $i$, including negative exponents.
\end{proof}

\paragraph{Exact scope.}
The formal module proves the displayed integer identities, product
divisibility, kernel lemma and finite-capacity implication. In the JP
application, actual modular reductions, the $n$-torsion sizes, product
height bounds and separation of distinct signed vectors are supplied
by the separate ordinary argument. The capacity theorem does not
construct those objects by taking them as hypotheses. Moreover,
$n^{|S|}$ cannot be replaced by $n$ merely because every target has
exponent dividing $n$. The actual relation $d=x+y$ is also a warning:
postulating an unproved radical bound for these three determinants
would introduce another ABC-shaped estimate.
